% Abhakseite „Das kann ich“ – Zone nur im Gesamt (\mitzone)
%% TODO Ich-kann-Sätze: Platzhalter wie in den Aufgabendateien
\begin{abhakseite}
\ifdefined\mitzone
\abhakgruppe{Kennst du schon}
\abhak{1}{Quadratische Gleichungen}
\abhak{2}{Äquivalenzumformungen und Probe}
\abhak{3}{Terme ausmultiplizieren, zusammenfassen, ausklammern, binomische Formeln vorwärts und rückwärts, Polynomdivision}
\abhak{4}{Potenzen, n-te Wurzeln und Logarithmus}
\abhak{5}{Die e-Funktion ist immer positiv und nie null, e hoch null ist eins}
\abhak{6}{Ableitung einer ganzrationalen Funktion und eines Produkts mit e-Funktion bilden}
\abhak{7}{Sinus am Einheitskreis}
\abhak{8}{Punkte am Graphen ablesen und zwei Graphenpunkte mit vorgegebenem Abstand in x und y finden, Maßstab (Längeneinheit in Meter)}
\abhak{9}{Äquivalenzumformungen und Probe – Fehler finden}
\abhak{10}{Äquivalenzumformungen und Probe – selbst rechnen}
\fi
\abhakgruppe{Einheit 1 · Ganzrationale Gleichungen}
\abhak{11}{Schnittpunkte und Stellen}
\abhak{12}{Waagerechte Ausdehnung eines Profils in vorgegebener Höhe über die Stellen mit vorgegebenem Funktionswert berechnen}
\abhak{13}{Schnittpunkte und Stellen – Fehler finden, Begründen, Anwendung, Darstellungswechsel}
\abhakgruppe{Einheit 2 · Gleichungen mit e-Funktion und Logarithmus}
\abhak{14}{Gleichungen mit e-Funktion und Logarithmus}
\abhak{15}{Schnittpunkte zweier Graphen durch Ausklammern eines gemeinsamen Faktors berechnen}
\abhak{16}{Gleichungen mit e-Funktion und Logarithmus – Fehler finden, Begründen, Anwendung}
\abhakgruppe{Einheit 3 · Gleichungen aufstellen und grafisch oder numerisch lösen}
\abhak{17}{Gleichungen aufstellen und grafisch oder numerisch lösen}
\abhak{18}{Schnittpunkte zweier Graphen mit dem Rechner ermitteln}
\abhak{19}{Gleichungen aufstellen und grafisch oder numerisch lösen – Fehler finden, Begründen, Anwendung, Darstellungswechsel}
\abhakgruppe{Einheit 4 · Ungleichungen}
\abhak{20}{Ungleichungen}
\abhak{21}{Intervall, in dem eine Modellfunktion mindestens einen vorgegebenen Wert annimmt, über eine Ungleichung bestimmen}
\abhak{22}{Ungleichungen – Fehler finden, Begründen, Anwendung, Darstellungswechsel}
\end{abhakseite}
